% Consolidated 24-electric-oscillations from 0d2a41587ad18cded191ec1d80e8718df37ec2fe; original TeX fragments preserved.

% === BEGIN 24-electric-oscillations-094-109.tex blob c96012258b552c75a97141e906a22ab0c0e0bb25 ===
% Irodov Part IV §4.2 94--109. Source user PDF pp.176--178;
% printed pp.167--169, source answers pp.317--318.
\electricoscproblem{94}{167}
铜金属薄板内自由电子整体沿板面法线微小平移 $x$ 后，表面产生电荷并形成静电回复场，引起电子等离子体振荡。铜中自由电子数密度 $n=0.85\times10^{29}\,\mathrm{m^{-3}}$，求其固有角频率。
\exsolution
离子晶格固定。电子整体位移 $x$ 在两表面形成电荷面密度大小 $\sigma=ne|x|$，内部净电场为 $E=\sigma/\varepsilon_0$，方向使电子回复。
单电子方程
$m_e\ddot x=-eE=-ne^2x/\varepsilon_0$，故
\[
\boxed{\omega_p=\sqrt{\frac{ne^2}{\varepsilon_0m_e}}
\simeq1.65\times10^{16}\,\mathrm{s^{-1}}}.
\]
\excheck 数密度单位必须为 $\mathrm{m^{-3}}$；该版底本纯文本把它识别为 $\mathrm{m^{-1}}$，与数密度定义不符。

\electricoscproblem{95}{168}
无电阻 LC 电路电容为 $C$、线圈自感为 $L$，持续自由无阻尼电振荡，电容两端电压最大幅度为 $V_m$。求任一时刻回路电流 $I$ 与电容电压 $V$ 的关系，分别采用电路方程和能量守恒推导。
\exsolution
回路方程 $L\dot I+V=0$，并有 $I=\pm C\dot V$（取一致的方向）。由消元得
$\ddot V+\omega_0^2V=0$，$\omega_0=1/\sqrt{LC}$。
若 $V=V_m\cos(\omega_0t+\alpha)$，
则 $I$ 与 $V$ 相差四分之一周期，
$I=\pm CV_m\omega_0\sin(\omega_0t+\alpha)$。
消去相位：
\[
\boxed{V^2+\frac LC I^2=V_m^2}.
\]
另由守恒
$\tfrac12CV^2+\tfrac12LI^2=\tfrac12CV_m^2$，也得到同一椭圆相轨迹。
\excheck $|I|\le V_m\sqrt{C/L}$；电容电压达到极值时电流恰为零。

\electricoscproblem{96}{168}
无电阻 LC 电路电容为 $C$、线圈自感为 $L$，初态开关断开，电容充至 $V_m$；$t=0$ 合闸。求：(a) $I(t)$；(b) 电容电场能与线圈磁场能相等瞬间的自感电动势。
\exsolution
设 $\omega_0=(LC)^{-1/2}$，
初态 $q(0)=CV_m$、$I(0)=0$。
电荷 $q=CV_m\cos\omega_0t$，
选放电方向为正，则
\[
\boxed{I(t)=V_m\sqrt{\frac CL}\sin\omega_0t}.
\]
电场与磁场能相等：
$CV^2/2=LI^2/2$，各占总能的一半，
故 $|V|=V_m/\sqrt2$。
依据无电阻 KVL，自感电动势与电容电压等模，
\[
\boxed{|\mathcal E_L|=\frac{V_m}{\sqrt2}}.
\]
\excheck 必须带 $1/\sqrt2$；该版书末公式文字抽取容易漏根号。

\electricoscproblem{97}{168}
无电阻 LC 振荡回路中，平行板电容器两极板被\emph{缓慢}拉开，使振荡固有频率提高至原来的 $\eta$ 倍；初始储能为 $W$。求移动极板所做功。
\exsolution
若极板分离远慢于振荡周期，系统为缓慢变参数谐振子。电容降低使
$\omega=(LC)^{-1/2}$ 增大。
对无耗散缓慢谐振，作用量
$J=W/\omega$ 为绝热不变量，所以
$W_f/W_i=\omega_f/\omega_i=\eta$，
由能量守恒，移动极板做功
\[
\boxed{A_{\rm adiabatic}=(\eta-1)W}.
\]
\textbf{书末条件差异：}该版书末给
$A=(\eta^2-1)W$。若额外要求在移动过程中电容电荷\emph{振幅固定}，则
$W\propto1/C\propto\omega^2$，确实出现该式；但无阻尼自由 LC 系统缓慢改变参数时电荷振幅并不固定。
\excheck 这两个结果属于\emph{不同的外加控制模型}。本稿保留原题“slowly”的通常绝热解释，并把书末差异列入勘误，而非无说明地替换。

\electricoscproblem{98}{168}
图4.27中电感 $L=2.5\,\mathrm{mH}$，两个电容 $C_1=2.0\,\mathrm{\mu F}$、
$C_2=3.0\,\mathrm{\mu F}$。两电容初始均充至同一电压 $V=180$V，合上中间电感支路开关后，求：(a) 固有振荡频率；(b) 电感电流最大值。
\exsolution
按原图两电容均跨接上下两公共节点，为并联，
$C_\Sigma=C_1+C_2=5.0\,\mathrm{\mu F}$。
总回路无电阻，固有角频率
$\omega_0=1/\sqrt{LC_\Sigma}$，
周期和普通频率为
\[
\boxed{T=2\pi\sqrt{LC_\Sigma}
\simeq0.70\,\mathrm{ms}},\qquad
\boxed{\nu_0=\frac1{2\pi\sqrt{LC_\Sigma}}
\simeq1.42\,\mathrm{kHz}}.
\]
初始电能全转为最大磁能：
$\tfrac12C_\Sigma V^2=\tfrac12LI_m^2$，
所以
\[
\boxed{I_m=V\sqrt{\frac{C_\Sigma}{L}}
\simeq8.0\,\mathrm A}.
\]
\excheck 书末(a)答的是周期 $T$，而题面只写“natural oscillation frequency”；本稿同时给出频率与周期，避免以毫秒冒充 Hz。

\electricoscproblem{99}{168}
图4.28的回路含两只相同电容 $C$、一只电感 $L$ 和开关，忽略电阻。左电容初始电压 $V_0$，右电容初始不带电；$t=0$ 合闸。求两电容各自电压随时间变化。
\exsolution
两电容经上方电感和下方导线耦合；以图中相同的上下极板参考方向记电压为 $V_1,V_2$。
两个电容上的总电荷守恒
$C(V_1+V_2)=CV_0$。
电感回路满足
$L\dot I=V_1-V_2$，并且
$I=-C\dot V_1=C\dot V_2$（选定方向可使两式同步）。
对差量 $u=V_1-V_2$，
$\ddot u+[2/(LC)]u=0$，
初值 $u(0)=V_0$、$\dot u(0)=0$。
故
\[
\boxed{V_1(t)=\frac{V_0}{2}(1+\cos\Omega t)},\quad
\boxed{V_2(t)=\frac{V_0}{2}(1-\cos\Omega t)},\quad
\boxed{\Omega=\sqrt{\frac2{LC}}}.
\]
\excheck $V_1+V_2=V_0$，初态分别为 $V_0,0$，半周期后两电容互换充电。

\electricoscproblem{100}{168}
无电阻 LC 电路电感为 $L$、电容为 $C$；线圈原处于恒定外磁场，穿过全部线圈的总外磁链为 $\Phi$。$t=0$ 将外磁场以远短于固有周期的时间撤去，求以后电路电流 $I(t)$。
\exsolution
外磁链骤降 $\Delta\Phi_{\rm ext}=-\Phi$，
根据短时间积分的法拉第—电感方程
$L\Delta I=-\Delta\Phi_{\rm ext}=\Phi$，
故撤场结束后初始电流
$I(0^+)=\Phi/L$。
撤场之后为自由 LC 振荡，
满足
$\ddot I+\omega_0^2I=0$，
并且电容原无电压，故
$\dot I(0^+)=0$。
于是
\[
\boxed{I(t)=\frac{\Phi}{L}
\cos\left(\frac{t}{\sqrt{LC}}\right)}.
\]
\excheck 磁链而不是单匝磁通入式；
撤场时间极短使电容电荷在该阶段几乎不变。

\electricoscproblem{101}{169}
电容电压按
$V(t)=V_me^{-\beta t}\cos\omega t$
作衰减振荡。求电压模在何时：(a) 与包络幅度相等；(b) 达到真正的局部极值。
\exsolution
(a) 电压模等于当前包络
$V_me^{-\beta t}$ 当且仅当
$|\cos\omega t|=1$，
所以
\[
\boxed{t_n^{\rm env}=\frac{n\pi}{\omega}},
\quad n=0,1,2,\ldots
\]
(b) 局部极值需 $\dot V=0$；
求导
$\dot V=-V_me^{-\beta t}
(\beta\cos\omega t+\omega\sin\omega t)$，
因此
\[
\boxed{\tan\omega t_n=-\frac\beta\omega},
\qquad
\boxed{t_n^{\rm ext}=
\frac{n\pi-\arctan(\beta/\omega)}{\omega}},
\]
选使 $t_n\ge0$ 的整数 $n$。
\excheck 因包络随时间衰减，电压实际局部极值一般不在包络接触点；初始 $t=0$ 是定义域端点。

\electricoscproblem{102}{169}
电阻为 $R$、电感为 $L$、电容为 $C$ 的串联 RLC 回路开关原断开，电容带初始电压 $V_0$，$t=0$ 合闸，开始欠阻尼振荡。求刚合闸时实际电压与位移振荡包络峰值之比。
\exsolution
令 $\beta=R/(2L)$、
$\omega_0=1/\sqrt{LC}$、
$\omega=\sqrt{\omega_0^2-\beta^2}$。
初始电流为零，即 $\dot q(0)=0$。
因此
$q(t)=CV_0e^{-\beta t}
[\cos\omega t+(\beta/\omega)\sin\omega t]$。
括号的相位合成振幅为
$\sqrt{1+(\beta/\omega)^2}=\omega_0/\omega$，
电容电压的指数包络峰值为
$V_m=V_0\omega_0/\omega$。
故
\[
\boxed{\frac{V_0}{V_m}
=\frac{\omega}{\omega_0}
=\sqrt{1-\frac{R^2C}{4L}}}.
\]
\excheck 欠阻尼条件为 $R^2C<4L$，根号非负；不得漏掉平方根。

\electricoscproblem{103}{169}
串联 RLC 回路电感 $L$、电容 $C$，自发自由阻尼电流
$I(t)=I_me^{-\beta t}\sin\omega t$。求电容电压 $V_C(t)$，特别是 $t=0$ 的电压。
\exsolution
回路方程
$L\dot I+RI+V_C=0$。
自由阻尼参数
$R=2L\beta$、
$\omega_0^2=\omega^2+\beta^2=1/(LC)$。
代入 $I(t)$，选电流与电容端电压的对应有向关系：
\[
\boxed{V_C(t)=
-LI_me^{-\beta t}
[\omega\cos\omega t+\beta\sin\omega t]}.
\]
写相位式，其振幅
$L I_m\sqrt{\omega^2+\beta^2}
=I_m\sqrt{L/C}$；
初始电压
\[
\boxed{|V_C(0)|=LI_m\omega
=I_m\sqrt{\frac{L/C}{1+\beta^2/\omega^2}}}.
\]
\excheck 不同电流箭头定义会使 $V_C$ 整体反号，但电压模与相位差保持不变。

\electricoscproblem{104}{169}
串联 RLC 电路 $C=4.0\,\mathrm{\mu F}$、
$L=2.0\,\mathrm{mH}$、$R=10\,\Omega$。
在自由阻尼振荡的某个电流极大时刻，求磁场能与电容电场能之比。
\exsolution
电流取极大瞬间 $\dot I=0$，
回路 KVL
$L\dot I+RI+V_C=0$
简化为
$V_C=-RI$。
故
\[
\frac{W_L}{W_C}
=\frac{LI^2/2}{CV_C^2/2}
=\frac{L}{CR^2}.
\]
代入得
\[
\boxed{\frac{W_L}{W_C}
=\frac{2.0\times10^{-3}}
{(4.0\times10^{-6})(10)^2}=5.0}.
\]
\excheck 电流最大时 $\dot I=0$，但\emph{电容电压并非零}；阻尼回路与无损 LC 的能量关系不同。

\electricoscproblem{105}{169}
串联 LC 振荡电路中，两只线圈的电感分别为 $L_1,L_2$、电阻分别为 $R_1,R_2$，忽略相互电感。希望用一只线圈替代它们，使振动频率与品质因数不变。求新线圈参数。
\exsolution
串联相同电流通过两线圈，磁链相加，
$L_{\rm eq}=L_1+L_2$，电阻亦串联相加
$R_{\rm eq}=R_1+R_2$。
电容 $C$ 保持不变时，
自由振动角频率
$\omega_d^2=1/(LC)-(R/2L)^2$，
品质因数由 $L,R,C$ 决定；
采用等效参数自然使两个量相同。
\[
\boxed{L=L_1+L_2,\qquad R=R_1+R_2}.
\]
\excheck 假设两线圈互感可忽略，
否则串联同向或反向耦合还需 $\pm2M$。

\electricoscproblem{106}{169}
LC 阻尼电路品质因数 $Q=5000$，振荡频率
$\nu=2.2\,\mathrm{MHz}$。求电流\emph{振幅}下降 $\eta=2.0$ 倍所需时间。
\exsolution
弱阻尼自由电流幅度
$I_m(t)=I_m(0)e^{-\beta t}$，
且 $Q\simeq\omega/(2\beta)$、
$\omega=2\pi\nu$，
所以
$\beta\simeq\pi\nu/Q$。
于是
\[
\boxed{t=\frac{\ln\eta}{\beta}
\simeq\frac{Q\ln\eta}{\pi\nu}}.
\]
代数值
\[
\boxed{t\simeq\frac{5000\ln2}{\pi(2.2\times10^6)}
=5.0\times10^{-4}\,\mathrm s=0.50\,\mathrm{ms}}.
\]
\excheck 该版书末的公式符合此计算，但数值提取似为“0.5 s”而非毫秒；代题面 MHz 时必为约0.5 ms。

\electricoscproblem{107}{169}
RLC 回路 $C=10\,\mathrm{\mu F}$、
$L=25\,\mathrm{mH}$、$R=1.0\,\Omega$。
求电流振幅降低 $e$ 倍需经历多少个\emph{实际阻尼振荡周期}。
\exsolution
$\beta=R/(2L)$、
$\omega_d=\sqrt{1/(LC)-\beta^2}$。
幅度降低 $e$ 倍需 $t=1/\beta=2L/R$，
期间周期数
\[
N=\frac t{2\pi/\omega_d}
=\frac{\omega_d}{2\pi\beta}
=\frac1{2\pi}\sqrt{\frac{4L}{CR^2}-1}.
\]
代入
\[
\boxed{N\simeq
\frac{\sqrt{4(0.025)/(10^{-5})-1}}{2\pi}
\simeq16}.
\]
\excheck 此处 $N$ 是周期数，不是电荷的极值次数，不能额外乘2。

\electricoscproblem{108}{169}
RLC 自由阻尼振荡电路品质因数为 $Q=5.0$，
求实际振荡角频率 $\omega$ 与无阻尼固有角频率 $\omega_0$ 的相对差百分比。
\exsolution
若品质因数按阻尼周期的 $\omega/(2\beta)$ 定义，
$Q=\omega/(2\beta)$，
则
$\omega_0^2=\omega^2+\beta^2
=\omega^2(1+1/(4Q^2))$，
因此
\[
\boxed{\frac{\omega_0-\omega}{\omega_0}
=1-\frac1{\sqrt{1+1/(4Q^2)}}}.
\]
$Q=5$ 时约
$\boxed{0.50\%}$。
\excheck 使用 $Q=\omega_0/(2\beta)$ 的另一常用定义也可在高 Q 下得到约0.5%，区别在更高阶修正。

\electricoscproblem{109}{169}
图4.29：电源电动势 $\mathcal E=2.0$V，
内阻 $r=9.0\,\Omega$，电容
$C=10\,\mathrm{\mu F}$，
电感 $L=100\,\mathrm{mH}$，
线圈电阻 $R=1.0\,\Omega$。
原先开关闭合达稳态，某刻断开，电容—线圈回路自由振荡。求断开后：(a) 初始振荡能量；(b) $t=0.30$s 的剩余能量。
\exsolution
闭合稳态电感短路，初始电流
$I_0=\mathcal E/(r+R)$，
电容电压 $V_{C0}=RI_0$。
因电容与线圈在稳态时为并联，
初始振荡能量
\[
W_0=\frac12LI_0^2+\frac12CV_{C0}^2
=\frac{\mathcal E^2(L+CR^2)}{2(r+R)^2}.
\]
数值
$\boxed{W_0\simeq2.0\,\mathrm{mJ}}$。
弱阻尼下能量包络衰减
$W(t)\simeq W_0e^{-Rt/L}$，
所以
\[
\boxed{W(0.30\,\mathrm s)\simeq
(2.0\,\mathrm{mJ})e^{-3}
\simeq0.10\,\mathrm{mJ}}.
\]
\excheck $r$ 在开关断开后被排除在振荡回路外，后续衰减取线圈电阻 $R$；能量指数衰减常数为 $R/L$ 而不是电流幅度的 $R/(2L)$。
% === END 24-electric-oscillations-094-109.tex ===

% === BEGIN 24-electric-oscillations-110-125.tex blob 9867215b006b63d31242264f054a194cc40b4fff ===
% Irodov Part IV §4.2 independent original 110--125.
% User PDF pp.179--181, printed pp.170--172, answers pp.318--319.
\electricoscproblem{110}{170}
某阻尼电振荡回路品质因数为 $Q=50$、无阻尼固有振荡频率
$\nu_0=5.5\,\mathrm{kHz}$。
问储存的振荡能量下降 $\eta=2.0$ 倍需要多久？
\exsolution
弱阻尼时电流、电压振幅随 $e^{-\beta t}$ 衰减，能量包络随 $e^{-2\beta t}$ 衰减。
品质因数 $Q\simeq\omega_0/(2\beta)$，
$\omega_0=2\pi\nu_0$，
所以能量下降倍数
$\eta=e^{2\beta t}=e^{\omega_0t/Q}$。
故
\[
\boxed{t\simeq\frac{Q\ln\eta}{2\pi\nu_0}
=\frac{50\ln2}{2\pi(5500)}
\simeq1.0\,\mathrm{ms}}.
\]
\excheck 能量衰减指数是\emph{振幅}衰减指数的两倍；若忽略该因子，结果会多一倍。

\electricoscproblem{111}{170}
电振荡电路含一只漏电的电容器，电容为 $C$，其漏电电阻为 $R$；与之连接的线圈自感为 $L$，线圈及导线电阻可忽略。
求：(a) 阻尼角频率；(b) 品质因数。
\exsolution
设电容两端电压为 $V$。电感支路电流为 $I_L$，
漏电流 $I_R=V/R$，
节点 KCL
$I_L+V/R+C\dot V=0$，
电感满足 $V=L\dot I_L$（取一致有向端子）。
对 KCL 求导并消去 $I_L$：
\[
\ddot V+\frac1{RC}\dot V+\frac1{LC}V=0.
\]
因此
\[
\boxed{\beta=\frac1{2RC}},\qquad
\boxed{\omega_d=\sqrt{\frac1{LC}-\frac1{4R^2C^2}}}.
\]
按 $Q=\omega_d/(2\beta)$，
\[
\boxed{Q=RC\omega_d
=\frac12\sqrt{\frac{4R^2C}{L}-1}}.
\]
\excheck 要有阻尼振荡必须 $4R^2C>L$；这个 $R$ 是\emph{跨电容漏电支路}，不同于串联 RLC 的线圈电阻。

\electricoscproblem{112}{170}
振荡回路 $C=2.0\,\mathrm{\mu F}$、
$L=5.0\,\mathrm{mH}$，
为维持电容电压幅度 $V_m=1.0$V 的无衰减振荡，平均需输入功率
$P=0.10\,\mathrm{mW}$。阻尼很小，求品质因数。
\exsolution
回路最大存储能量
$W=\tfrac12CV_m^2$，
每周损失能量约 $P(2\pi/\omega_0)$，
$Q=2\pi W/(\text{每周耗能})=\omega_0W/P$。
因 $\omega_0=(LC)^{-1/2}$，
\[
\boxed{Q=\frac{V_m^2}{2P}\sqrt{\frac CL}}
=\frac{1}{2(10^{-4})}\sqrt{\frac{2\times10^{-6}}{5\times10^{-3}}}
=\boxed{100}.
\]
\excheck $P$ 必须以瓦特而不是毫瓦特代入；
$Q$ 为无量纲。

\electricoscproblem{113}{170}
电振荡回路中的有功电阻 $R=0.45\,\Omega$，
要维持电流幅度
$I_m=30\,\mathrm{mA}$ 的无衰减正弦振荡，
需要输入多大平均功率？
\exsolution
电流 $I=I_m\cos\omega t$，
焦耳热平均功率
$\overline P=R\langle I^2\rangle$
且 $\langle\cos^2\rangle=1/2$，
\[
\boxed{\overline P=\frac12RI_m^2}
=\frac12(0.45)(0.030)^2\,\mathrm W
=\boxed{2.0\times10^{-4}\,\mathrm W
=0.20\,\mathrm{mW}}.
\]
\excheck 书后答案将相同解析公式后标成“20 mW”，与 $30$mA、$0.45\Omega$ 直接代入差100倍。保持原版数据并登记勘误。

\electricoscproblem{114}{170}
振荡回路 $C=1.2\,\mathrm{nF}$、
$L=6.0\,\mathrm{\mu H}$、$R=0.50\,\Omega$，
要维持电容电压幅度
$V_m=10$V 的无衰减振荡，所需平均输入功率为多少？
\exsolution
弱阻尼时电感电流幅度
$I_m=\omega_0CV_m=
V_m\sqrt{C/L}$，
损耗平均值
$\overline P=\tfrac12 R I_m^2$。
因此
\[
\boxed{\overline P=\frac{RCV_m^2}{2L}}
=\frac{(0.50)(1.2\times10^{-9})(10)^2}
{2(6.0\times10^{-6})}
=\boxed{5.0\,\mathrm{mW}}.
\]
\excheck 这里 $V_m$ 是振幅；不必先换算有效值，只需显式保留因子1/2。

\electricoscproblem{115}{170}
图4.30：电容 $C$、线圈自感 $L$ 以及电阻 $R$ 分别连接在同一上、下两端点之间（并联）。
求自由阻尼振荡角频率，以及发生振荡对 $C,L,R$ 的约束。
\exsolution
图中为\emph{并联}的 L、C、R 三支路。
令两端电压为 $V$，
节点电流守恒
$C\dot V+V/R+I_L=0$。
电感关系 $L\dot I_L=V$；
对 KCL 求导：
\[
\ddot V+\frac1{RC}\dot V+\frac1{LC}V=0.
\]
因此
\[
\boxed{\omega_d=
\sqrt{\frac1{LC}-\frac1{4R^2C^2}}},
\]
仅当
\[
\boxed{R>\frac12\sqrt{\frac LC}}
\]
时有欠阻尼振荡。
\excheck 书后对并联电阻的门限是\emph{大电阻}更容易振荡；不能错套串联电阻 $R<2\sqrt{L/C}$ 的门限。

\electricoscproblem{116}{170}
图4.31有两种电振荡线路，电容均为 $C$。
(a) 两只各具自感 $L_1,L_2$、电阻 $R_1,R_2$ 的线圈与电容并联，线圈之间互感忽略；
(b) 一只自感 $L$、电阻 $R$ 的线圈与电容并联。
求使两种自由振荡的频率与阻尼相同所需的参数关系。
\exsolution
对(a)中总电感支路，在\emph{低损耗或两支路阻尼比率相容}时，可用并联电感与阻尼等效模型表示。
两自感的并联磁能等效给
\[
\boxed{\frac1L=\frac1{L_1}+\frac1{L_2}}.
\]
书后采用并联损耗的等效电导加法：
\[
\boxed{\frac1R=\frac1{R_1}+\frac1{R_2}}.
\]
在该等效模型下，
$\omega_0=1/\sqrt{LC}$、
$\beta=1/(2RC)$ 相同，因而阻尼自由频率相同。
\excheck 对图(a)中两条真正的\emph{串联 RL 支路}，
其总导纳是
$(R_1+sL_1)^{-1}+(R_2+sL_2)^{-1}$，
并不对任意参数严格等于
$(R+sL)^{-1}$。
若要求完整暂态严格等效，还需对各支路的 $R_i/L_i$ 施加匹配条件；本题书末的两条倒数关系只能在其简化电路模型下直接使用。

\electricoscproblem{117}{170}
串联 RLC 回路电容原先充电至 $V_0$，
合闸瞬间取 $t=0$。
回路电阻正好为临界阻尼值。
求随后电流 $I(t)$ 和峰值 $I_{\max}$。
\exsolution
临界阻尼
$R_{\rm cr}=2\sqrt{L/C}$，
$\beta=R/(2L)=1/\sqrt{LC}$。
电荷满足
$\ddot q+2\beta\dot q+\beta^2q=0$，
初始 $q(0)=CV_0$、$\dot q(0)=0$，
故
$q(t)=CV_0(1+\beta t)e^{-\beta t}$。
放电电流
$I=-\dot q=CV_0\beta^2 t\,e^{-\beta t}$，即
\[
\boxed{I(t)=\frac{V_0}{L}\,t\,
e^{-t/\sqrt{LC}}}.
\]
极大点 $t_m=1/\beta=\sqrt{LC}$，
\[
\boxed{I_{\max}=\frac{V_0}{e}\sqrt{\frac CL}}.
\]
\excheck 初始电流为零，$I$ 先增后减，始终不发生往复振荡，符合临界阻尼性质。

\electricoscproblem{118}{171}
电感 $L$、有功电阻 $R$ 的线圈在 $t=0$ 接入
$V(t)=V_m\cos\omega t$ 的正弦电压源。
起初无电流，求开通后 $I(t)$。
\exsolution
方程
$L\dot I+RI=V_m\cos\omega t$。
稳态项
$I_s=I_m\cos(\omega t-\varphi)$，
其中
$I_m=V_m/\sqrt{R^2+\omega^2L^2}$、
$\tan\varphi=\omega L/R$。
齐次项 $I_h=Ae^{-Rt/L}$；
由 $I(0)=0$ 有
$A=-I_m\cos\varphi$。
所以
\[
\boxed{I(t)=\frac{V_m}{\sqrt{R^2+\omega^2L^2}}
\left[\cos(\omega t-\varphi)
-\cos\varphi\,e^{-Rt/L}\right]},
\quad
\varphi=\arctan\frac{\omega L}{R}.
\]
\excheck $I(0)=0$；电感导致相位\emph{滞后}而非超前。

\electricoscproblem{119}{171}
电容 $C$ 与电阻 $R$ 串联，在 $t=0$ 接入
$V(t)=V_m\cos\omega t$ 的电源。
求开通后回路电流 $I(t)$，
并区分不同初始电容电压条件。
\exsolution
电路方程
$RI+q/C=V_m\cos\omega t$，
$I=\dot q$。
稳态电流
\[
I_s=I_m\cos(\omega t+\theta),\quad
I_m=\frac{V_m}{\sqrt{R^2+1/(\omega C)^2}},
\quad
\tan\theta=\frac1{\omega RC}.
\]
总电流
$I(t)=I_s+A e^{-t/(RC)}$，
由开通瞬间的电容电压 $V_C(0)=V_{C0}$，
$I(0)=(V_m-V_{C0})/R$，故
\[
\boxed{I(t)=I_m\cos(\omega t+\theta)
+\left[\frac{V_m-V_{C0}}R-I_m\cos\theta\right]
e^{-t/(RC)}}.
\]
若电容初始不带电 $V_{C0}=0$，直接代入即可。
\textbf{书末约定：}书末给出的形式
$I=I_m[\cos(\omega t-\varphi)-
\cos\varphi e^{-t/(RC)}]$（$\varphi=-\theta$）
对应 $I(0)=0$，即预先令
$V_{C0}=V_m$，而\emph{不是}初始未充电状态。
\excheck 原题未指定初始电容电压，故单一暂态解不唯一；本稿先给包含 $V_{C0}$ 的完整解，再说明书末隐含初值。

\electricoscproblem{120}{171}
长单层密绕螺线管绕线电阻率为 $\rho$，
单位轴向长度匝数 $n$，
螺线管截面半径为 $a$，
绝缘层厚度可忽略。
施加频率 $\nu$ 的交变电压时，
求电流与电压的相位差。
\exsolution
导线直径约 $d=1/n$，
截面积 $A_w=\pi/(4n^2)$。
每单位螺线管轴向长度的导线长度
$s'=2\pi an$，所以
$R'=\rho s'/A_w=8\rho an^3$。
长螺线管单位轴长自感
$L'=\mu_0n^2\pi a^2$。
交变角频率 $\Omega=2\pi\nu$，
相位滞后
$\tan\varphi=\Omega L'/R'$，
故
\[
\boxed{\tan\varphi=
\frac{\mu_0\pi^2\nu a}{4\rho n}}.
\]
\excheck 该表达式适用长、单层、紧密绕制和忽略趋肤效应；$n$ 是每米匝数，不能当总匝数。

\electricoscproblem{121}{171}
电容 $C$ 与电阻 $R=110\,\Omega$ 串联，
接到电压幅度 $V_m=110$V 的交流源，
测得稳态电流幅度 $I_m=0.50$A。
求电流与源电压的相位差。
\exsolution
阻抗模
$|Z|=V_m/I_m=220\,\Omega$。
串联 RC 中
$\cos\varphi=R/|Z|=110/220=1/2$。
因此电流相对电源电压\emph{超前}
\[
\boxed{\varphi=\arccos(1/2)=60^\circ}.
\]
\excheck 电路为容性，电流必须超前；若以电压相对电流的相位定义，则记为 $-60^\circ$。

\electricoscproblem{122}{171}
图4.32给出最简脉动滤波器：输入
$V_{\rm in}(t)=V_0(1+\cos\omega t)$，
经串联电阻 $R$ 后输出跨并联电容 $C$ 的电压。
求：(a) 稳态输出；(b) $\omega=314\,\mathrm{s^{-1}}$ 时，
输出交流幅度为输出直流分量的 $1/\eta$，
$\eta=7.0$，所需 $RC$。
\exsolution
直流稳态下电容开路，输出常值部分为 $V_0$。
交流部分采用传递函数
$H(i\omega)=1/(1+i\omega RC)$，
因此
\[
\boxed{V_{\rm out}(t)=V_0+
\frac{V_0}{\sqrt{1+(\omega RC)^2}}
\cos\left(\omega t-\arctan\omega RC\right)}.
\]
要求
$\sqrt{1+(\omega RC)^2}=\eta$，
所以
\[
\boxed{RC=\frac{\sqrt{\eta^2-1}}{\omega}
=\frac{\sqrt{48}}{314}
\simeq22\,\mathrm{ms}}.
\]
\excheck 输出直流增益为1，交流增益小于1；频率增大则滤波效果增强。

\electricoscproblem{123}{171}
图4.33(a)(b)给出两种正弦交流串联电路，
要求按外加电压与电流的相位关系作近似电压相量图。
(a) 纯电阻 $R_0$ 与具有自感 $L$、电阻 $R$ 的线圈串联；
(b) 电容 $C$ 与具有自感 $L$、电阻 $R$ 的线圈串联，
且题设 $\omega<\omega_0=1/\sqrt{LC}$。
\exsolution
取电流相量 $\bm I$ 沿实轴。
(a) 电阻 $R_0$ 上电压
$\underline V_{R_0}=IR_0$ 同相，
线圈电压
$\underline V_{\rm coil}=I(R+i\omega L)$，
总电压
\[
\boxed{\underline V=I[(R_0+R)+i\omega L]},
\quad
\tan\varphi=\frac{\omega L}{R_0+R}>0.
\]
相量图中两段水平电阻电压相加，感抗电压竖直向上，总电压在第一象限。
(b) 电容电压
$\underline V_C=-iI/(\omega C)$，
线圈电压如上；
$\omega<\omega_0$ 时
$\omega L-1/(\omega C)<0$，故
\[
\boxed{\underline V=I\left[
R+i\left(\omega L-\frac1{\omega C}\right)\right]},
\]
总电压相对电流落在第四象限；相量图中的电容电压向下长于电感向上的电压。
\excheck 用复数相量的实轴和正负虚轴描述给出了原图4.33的两个电压矢量图的全部相对方向与长度约束；不用电压有效值和峰值混算。

\electricoscproblem{124}{171}
串联 RLC 回路含
$C=22\,\mathrm{\mu F}$、线圈
$R=20\,\Omega$、$L=0.35$H，
外加电压幅度 $V_m=180$V，
角频率 $\omega=314\,\mathrm{s^{-1}}$。
求：(a) 电流幅度；(b) 电流与外电压相位差；
(c) 电容及线圈\emph{整个线圈端子}的电压幅度。
\exsolution
感抗 $X_L=\omega L\simeq109.9\,\Omega$，
容抗 $X_C=1/(\omega C)\simeq144.8\,\Omega$。
阻抗模
$Z=\sqrt{R^2+(X_L-X_C)^2}\simeq40.2\,\Omega$。
故
\[
\boxed{I_m=\frac{V_m}{Z}\simeq4.5\,\mathrm A}.
\]
以源电压相对电流的相位为 $\varphi$：
\[
\boxed{\tan\varphi=\frac{X_L-X_C}{R}},
\quad\boxed{\varphi\simeq-60^\circ},
\]
即电流超前电压约 $60^\circ$。
电容两端幅度
$\boxed{V_{Cm}=I_mX_C\simeq0.65\,\mathrm{kV}}$；
整个线圈含内阻，其端电压幅度
\[
\boxed{V_{{\rm coil},m}
=I_m\sqrt{R^2+X_L^2}\simeq0.50\,\mathrm{kV}}.
\]
\excheck 交流相量叠加不等于电压幅值的代数相加，
因此单个器件的电压幅度可大于电源幅度。

\electricoscproblem{125}{171--172}
正弦电压幅度保持不变、频率可调的电源，
向电容 $C$、纯电阻 $R$、无电阻电感 $L$ 的串联回路供电。
求令电压\emph{幅度}最大的驱动频率：
(a) 跨电容；
(b) 跨电感。
\exsolution
记 $\omega_0^2=1/(LC)$、
$\beta=R/(2L)$，驱动频率为 $\Omega$。
电流幅度
$I_m=V_m/\sqrt{R^2+(\Omega L-1/(\Omega C))^2}$。
电容电压
$V_C=I_m/(\Omega C)$，
整理平方：
\[
V_C^2=\frac{V_m^2\omega_0^4}
{(\omega_0^2-\Omega^2)^2+4\beta^2\Omega^2}.
\]
分母对 $\Omega^2$ 求极小，
\[
\boxed{\Omega_{C,\max}=
\sqrt{\omega_0^2-2\beta^2}}.
\]
电感电压
$V_L=I_m\Omega L$，对应
\[
V_L^2=\frac{V_m^2\Omega^4}
{(\omega_0^2-\Omega^2)^2+4\beta^2\Omega^2}.
\]
令 $y=\Omega^2$，对上式求导得极值
\[
\boxed{\Omega_{L,\max}
=\frac{\omega_0^2}
{\sqrt{\omega_0^2-2\beta^2}}}.
\]
\excheck 两式要求 $\omega_0^2>2\beta^2$；
高阻尼时不存在同样的内部有限频率幅度峰。
% === END 24-electric-oscillations-110-125.tex ===

% === BEGIN 24-electric-oscillations-126-137.tex blob 96bab7cd303d468988481af6ba4ab13ad8be6ea5 ===
% Irodov Part IV §4.2 original 126--137. Source PDF pp.181--182,
% printed pp.172--173. Answers PDF pp.319--320.
\electricoscproblem{126}{172}
电感 $L=0.36\,\mathrm H$、有功电阻
$R=40\,\Omega$ 的线圈与电容 $C$ 串联，
外加正弦电压角频率 $\omega=314\,\mathrm{s^{-1}}$、
幅度 $V_m=180$V。
调节 $C$ 使\emph{整个线圈}端电压幅度最大；
求最佳电容以及对应线圈、电容电压幅度。
\exsolution
线圈幅度阻抗
$|Z_L|=\sqrt{R^2+(\omega L)^2}$
在调节 $C$ 时不变；线圈端电压
$V_L=V_m|Z_L|/
\sqrt{R^2+(\omega L-1/(\omega C))^2}$。
分母在总反应抗抵消时最小：
$1/(\omega C)=\omega L$，
故
\[
\boxed{C_*=\frac1{\omega^2L}
\simeq28\,\mathrm{\mu F}}.
\]
此时 $I_m=V_m/R$，
\[
\boxed{V_{{\rm coil},m}=
V_m\sqrt{1+\left(\frac{\omega L}{R}\right)^2}
\simeq0.54\,\mathrm{kV}},
\quad
\boxed{V_{C,m}=\frac{V_m\omega L}{R}
\simeq0.51\,\mathrm{kV}}.
\]
\excheck 两器件端电压的幅度均可高于源幅度，源电压为它们的相量和。

\electricoscproblem{127}{172}
介质有微弱导电性的电容器，电容为 $C$、漏电电阻为 $R$，
两端接正弦电压
$V(t)=V_m\cos\omega t$。
导线电阻忽略，求稳态供电导线中的电流随时间的变化。
\exsolution
漏电传导电流
$I_R=V/R=(V_m/R)\cos\omega t$；
位移电流
$I_C=C\dot V=-\omega CV_m\sin\omega t$。
总供电电流
\[
I(t)=\frac{V_m}{R}\cos\omega t-\omega CV_m\sin\omega t.
\]
写成相位形式
\[
\boxed{I(t)=I_m\cos(\omega t+\varphi)},\qquad
\boxed{I_m=\frac{V_m}{R}\sqrt{1+(\omega RC)^2}},
\quad
\boxed{\tan\varphi=\omega RC}.
\]
\excheck 导线电流包括漏电传导与充放电电流；
它超前电压，且 $\omega\to0$ 回到纯电阻漏电流。

\electricoscproblem{128}{172}
LC 回路电容为 $C$、自感为 $L_1$；
其线圈与另一只理想短路线圈（自感为 $L_2$）互感为 $M$。
短路线圈有功电阻忽略。
求该 LC 回路的固有角频率。
\exsolution
两线圈的有向磁链
$\Psi_1=L_1I_1+MI_2$、
$\Psi_2=MI_1+L_2I_2$。
次级线圈短路且无电阻，
其磁链变化为零；若初始磁链为零，
$\Psi_2=0$，所以
$I_2=-(M/L_2)I_1$。
代回初级磁链
$\Psi_1=(L_1-M^2/L_2)I_1$，
等效自感
$L_{\rm eff}=L_1-M^2/L_2$。
故
\[
\boxed{\omega_0=
\frac1{\sqrt{C(L_1-M^2/L_2)}}
=\sqrt{\frac{L_2}{C(L_1L_2-M^2)}}}.
\]
\excheck 由磁能正定性 $M^2<L_1L_2$，
等效电感必须为正；电磁耦合缩小初级电感，从而提高频率。

\electricoscproblem{129}{172}
串联谐振回路接交流电源，在电流共振时
电容器两端电压幅度是外加电源幅度的 $n$ 倍。
求品质因数 $Q$。
\exsolution
串联电流共振频率
$\omega_0=1/\sqrt{LC}$，
电路阻抗最小为 $R$，
$I_m=V_m/R$。
电容电压幅度与源电压幅度之比
\[
n=\frac{I_m/(\omega_0 C)}{V_m}
=\frac1{\omega_0CR}
=\frac{\omega_0L}{R}.
\]
若品质因数按\emph{实际阻尼频率}
$\omega_d=\sqrt{\omega_0^2-\beta^2}$ 定义：
$Q=\omega_d/(2\beta)$，
$\beta=R/(2L)$，
则
\[
\boxed{Q=\sqrt{\left(\frac{\omega_0L}{R}\right)^2
-\frac14}
=\sqrt{n^2-\frac14}}.
\]
\excheck 弱阻尼时 $Q\simeq n$；区别来自采用 $\omega_d$ 还是 $\omega_0$ 定义 Q，而不是测量电压出现错误。

\electricoscproblem{130}{172}
串联 RLC 回路的电源角频率与电压幅度保持不变，
调整自感 $L$ 使电流达到最大幅度。
若此时自感增加到原来的 $n$ 倍，电流幅度降为原来的 $1/\eta$，
求原共振状态的品质因数。
\exsolution
调节 $L$ 使
$\omega L=1/(\omega C)$，
原最大电流 $I_m=V_m/R$。
自感换为 $nL$ 后，
反应抗
$X'=\omega(nL)-1/(\omega C)
=\omega L(n-1)$，
所以
\[
\eta^2=
\frac{R^2+[\omega L(n-1)]^2}{R^2}
=1+\left(\frac{\omega L}{R}\right)^2(n-1)^2.
\]
由第129题的品质因数约定
$Q^2=(\omega L/R)^2-1/4$，
得
\[
\boxed{Q=
\sqrt{\frac{\eta^2-1}{(n-1)^2}-\frac14}}.
\]
\excheck 需有 $n\ne1$、$\eta\ge1$；
此关系还给出数据相容性 $Q^2\ge0$。

\electricoscproblem{131}{172}
串联 RLC 回路交流源电压幅度恒定，
改变电源频率时，在
$\omega_1,\omega_2$ 两处电流幅度都比电流共振峰小 $n$ 倍。
求：(a) 共振角频率；(b) 品质因数。
\exsolution
电流幅度
$I(\omega)=V_m/\sqrt{R^2+(\omega L-1/(\omega C))^2}$。
在电流峰 $\omega_0^2=1/(LC)$，
$I_{\max}=V_m/R$。
半峰倍率条件
\[
\left(\frac{\omega_0^2-\omega_i^2}{\omega_i}\right)^2
=(n^2-1)\left(\frac RL\right)^2.
\]
较低和较高两正根对应相反符号，
解得
$\omega_1\omega_2=\omega_0^2$ 与
$\omega_2-\omega_1=(R/L)\sqrt{n^2-1}$。
故
\[
\boxed{\omega_0=\sqrt{\omega_1\omega_2}}.
\]
按 $Q=\omega_dL/R$，$\omega_d^2=\omega_0^2-(R/2L)^2$：
\[
\boxed{Q=
\sqrt{\frac{\omega_1\omega_2(n^2-1)}
{(\omega_2-\omega_1)^2}-\frac14}}.
\]
\excheck 这里 $\omega_1,\omega_2$ 必须位于峰的两侧且 $n>1$。

\electricoscproblem{132}{172}
证明弱阻尼串联谐振回路的品质因数
$Q\simeq\omega_0/\Delta\omega$，
其中 $\Delta\omega$ 是电流谐振曲线在
\emph{峰值幅度的 $1/\sqrt2$（半功率）}处的频宽。
\exsolution
在 $\omega$ 接近 $\omega_0$ 时，
串联反应抗
$X(\omega)=\omega L-1/(\omega C)
\simeq2L(\omega-\omega_0)$。
所以
\[
\frac{I(\omega)^2}{I_{\max}^2}
\simeq\frac{R^2}{R^2+
4L^2(\omega-\omega_0)^2}.
\]
半功率点
$I^2=I_{\max}^2/2$
对应
$|2L(\omega-\omega_0)|=R$，
两侧点间宽度
$\Delta\omega=R/L$。
弱阻尼
$Q\simeq\omega_0L/R$，
故
\[
\boxed{Q\simeq\frac{\omega_0}{\Delta\omega}}.
\]
\excheck 底本文字抽取把“$\sqrt2$ times less”中的根号漏掉，只剩“2 times less”；若按幅度真正降为1/2，则频宽系数不同，不能推导此式。

\electricoscproblem{133}{172--173}
串联 RLC 回路先后接入电压幅度相同、频率分别为
$\omega_0$ 和 $\eta\omega_0$ 的两个交流源，
其中 $\omega_0=1/\sqrt{LC}$。
若品质因数为 $Q$，
求两情况下电流幅度之比 $I_0/I$；
并在 $\eta=1.10$ 时分别算
$Q=10$ 和 $Q=100$。
\exsolution
在 $\omega_0$ 处阻抗是 $R$，
$I_0=V_m/R$。
在 $\eta\omega_0$ 处的反应抗
\[
X=\eta\omega_0L-\frac1{\eta\omega_0C}
=\omega_0L\frac{\eta^2-1}{\eta}.
\]
故
\[
\frac{I_0}{I}=
\sqrt{1+\left(\frac X R\right)^2}.
\]
用本书阻尼品质因数定义
$\omega_0^2L^2/R^2=Q^2+1/4$，
得
\[
\boxed{\frac{I_0}{I}=
\sqrt{1+\frac{(Q^2+1/4)(\eta^2-1)^2}{\eta^2}}}.
\]
代 $\eta=1.10$：
\[
\boxed{I_0/I\simeq2.2\quad(Q=10)},\qquad
\boxed{I_0/I\simeq19\quad(Q=100)}.
\]
\excheck 高品质因数回路对较小的失谐更敏感；
两个数值均符合书末。

\electricoscproblem{134}{173}
蓄电池以恒定直流电流 $I_0$ 充电需 $t_0$ 小时。
若改用半波整流交流电流充电，
整流电流的\emph{有效值}仍为 $I_0$，
问需要多长时间，忽略实际电池充电效率的变化。
\exsolution
半波整流正弦电流一个周期中
$I(t)=I_m\sin\omega t$（$0<\omega t<\pi$），
另一半周期为零。
有效值平方
$I_{\rm rms}^2=(2\pi)^{-1}\int_0^\pi I_m^2\sin^2u\,\mathrm du=I_m^2/4$，
故 $I_m=2I_0$。
周期平均值
$\overline I=(2\pi)^{-1}\int_0^\pi I_m\sin u\,\mathrm du=I_m/\pi=2I_0/\pi$。
相同总电荷
$I_0t_0=\overline It$，所以
\[
\boxed{t=\frac\pi2t_0}.
\]
\excheck 有效电流相同不等于平均输送电荷速率相同；半波电流的峰值与平均值都必须单独计算。

\electricoscproblem{135}{173}
一股周期电流在整个周期的平均值为 $I_0$。
求电流的有效值：
(a) 如原图4.34，电流在每一周期从零\emph{线性增加到峰值}后瞬间归零的锯齿波；
(b) $I(t)\propto|\sin\omega t|$。
\exsolution
(a) 令周期为 $T$，
$I(t)=I_m t/T$（$0<t<T$）。
平均
$I_0=T^{-1}\int_0^T I_m t/T\,\mathrm dt=I_m/2$；
有效值
\[
I_{\rm rms}^2=\frac1T\int_0^T
\left(\frac{I_mt}{T}\right)^2\mathrm dt
=\frac{I_m^2}{3}.
\]
因 $I_m=2I_0$，
\[
\boxed{I_{{\rm rms},a}=\frac2{\sqrt3}I_0\simeq1.15I_0}.
\]
(b) 对全波整流正弦
$\overline I=(2/\pi)I_m=I_0$，
$I_{\rm rms}=I_m/\sqrt2$，
故
\[
\boxed{I_{{\rm rms},b}=\frac{\pi}{2\sqrt2}I_0
\simeq1.11I_0}.
\]
\excheck (a) 的波形来自原图4.34，不能误认成三角上下对称波而改变积分。

\electricoscproblem{136}{173}
线圈 $L=7\,\mathrm{mH}$、
$R=44\,\Omega$，先接直流源 $V_0$，
再接有效电压同为 $V=V_0$ 的正弦交流源。
求交流频率 $\nu$，
使线圈交流有功热功率只有直流情形的 $1/\eta$，
其中 $\eta=5.0$。
\exsolution
直流稳定时 $P_{\rm DC}=V_0^2/R$；
交流时总阻抗
$Z=\sqrt{R^2+(2\pi\nu L)^2}$，
有功功率
$P_{\rm AC}=V_0^2R/Z^2$。
比值
$P_{\rm DC}/P_{\rm AC}
=Z^2/R^2=1+(2\pi\nu L/R)^2=\eta$，
所以
\[
\boxed{\nu=\frac{R}{2\pi L}\sqrt{\eta-1}
=\frac{44}{2\pi(0.007)}(2)
\simeq2.0\,\mathrm{kHz}}.
\]
\excheck 电压给的是有效值，因而功率表达式不用再额外除2。

\electricoscproblem{137}{173}
线圈感抗为 $X_L=30\,\Omega$，
总阻抗模为 $Z=50\,\Omega$，
接有效电压 $V=100$V 的交流电源。
求电流相对电压的相位差，以及线圈平均发热功率。
\exsolution
线圈阻抗为 $R+iX_L$，
$R=\sqrt{Z^2-X_L^2}
=\sqrt{2500-900}=40\,\Omega$。
相位滞后角
\[
\boxed{\varphi=\arctan\frac{X_L}{R}
=\arctan\frac{30}{40}
\simeq37^\circ}.
\]
电流有效值 $I=V/Z=2.0$A，
热功率
\[
\boxed{P=I^2R=(2.0)^2(40)
=160\,\mathrm W=0.16\,\mathrm{kW}}.
\]
\excheck 感性线圈中电流\emph{滞后}电压，只有其有功电阻产生焦耳热。
% === END 24-electric-oscillations-126-137.tex ===

% === BEGIN 24-electric-oscillations-138-149.tex blob aff6629b320b3a0f52111585878aca06c8bda618 ===
% Irodov Part IV §4.2 original 138--149 (end of electric oscillations).
% Checked user source PDF pp.182--184, printed pp.173--175 and figures 4.34--4.36.
\electricoscproblem{138}{173}
电感 $L=0.70\,\mathrm H$、线圈内阻 $r=20\,\Omega$ 与外接无感电阻 $R$ 串联，
加有效电压 $V=220$V、角频率
$\omega=314\,\mathrm{s^{-1}}$。
求使整个电路产生的热功率最大时 $R$ 的数值与最大功率。
\exsolution
总有功电阻
$u=R+r$，感抗 $X_L=\omega L$，
有功功率
$P(u)=V^2u/[u^2+X_L^2]$。
对 $u$ 求导：
\[
\frac{\mathrm dP}{\mathrm du}
=V^2\frac{X_L^2-u^2}{(u^2+X_L^2)^2}.
\]
极大值条件 $u=X_L$。
故
\[
\boxed{R_*=\omega L-r
=314(0.70)-20\simeq0.20\,\mathrm{k\Omega}},
\]
\[
\boxed{P_{\max}=\frac{V^2}{2\omega L}
\simeq0.11\,\mathrm{kW}}.
\]
\excheck 最优 $R\ge0$ 需 $\omega L\ge r$；本题满足。题目问\emph{整个串联回路}总热功率，不能只算外接电阻部分。

\electricoscproblem{139}{173}
串联的电容和有内阻的线圈接交流电网，
调节电容使线圈产生的热功率增加
$n=1.7$ 倍。求功率因数 $\cos\varphi$ 改变百分之多少。
\exsolution
外加电源有效电压 $V$ 固定，线圈内阻为 $R$，
串联阻抗
$Z=\sqrt{R^2+X^2}$。
热功率
$P=I^2R=V^2R/Z^2
=(V^2/R)\cos^2\varphi$，
其中 $\cos\varphi=R/Z$。
两状态
$P_2/P_1=n$
使
$(\cos\varphi_2/\cos\varphi_1)^2=n$，
因功率因数取非负值，
\[
\boxed{\frac{\cos\varphi_2-\cos\varphi_1}
{\cos\varphi_1}\,100\%
=(\sqrt n-1)\,100\%\simeq30\%}.
\]
\excheck 该结论只需电源电压与线圈电阻不变；无需已知具体电容值。

\electricoscproblem{140}{173}
品质因数 $Q=100$ 的串联振荡回路，
接有效电压恒定的正弦电源。
在共振频率处热功率最大；
求将频率相对共振频率偏移百分之多少，
使有功功率降低为原来的 $1/n$，
其中 $n=2.0$。
\exsolution
弱阻尼近共振条件下
$X=\omega L-1/(\omega C)
\simeq2L(\omega-\omega_0)$。
最大功率
$P_{\max}=V^2/R$，
失谐时
$P/P_{\max}=R^2/(R^2+X^2)$。
令比值为 $1/n$，
$|X|=R\sqrt{n-1}$。
以 $Q\simeq\omega_0L/R$，
\[
\boxed{\left|\frac{\Delta\omega}{\omega_0}\right|
\simeq\frac{\sqrt{n-1}}{2Q}
=\frac1{200}=0.005=0.50\%}.
\]
\excheck 此为高Q、频偏很小的近共振公式，适用于两侧对称的近似半功率点。

\electricoscproblem{141}{173}
无感电阻 $R=0.16\,\mathrm{k\Omega}$
与具有有功内阻的线圈串联，
接有效电压 $V=220$V。
测得无感电阻两端有效电压
$V_1=80$V、线圈两端有效电压
$V_2=180$V。
求线圈内部发热功率。
\exsolution
串联电流有效值
$I=V_1/R$。
因纯电阻的电压 $\underline V_1$ 与电流同相，
线圈电压 $\underline V_2$ 相对电流滞后（电压超前） $\varphi$，
总相量
$\underline V=\underline V_1+\underline V_2$。
因此
$V^2=V_1^2+V_2^2+2V_1V_2\cos\varphi$，
线圈功率
$P_{\rm coil}=IV_2\cos\varphi
=(V_1/R)V_2\cos\varphi$。
消去 $\cos\varphi$：
\[
\boxed{P_{\rm coil}=
\frac{V^2-V_1^2-V_2^2}{2R}}
=\frac{220^2-80^2-180^2}{2(160)}
=\boxed{30\,\mathrm W}.
\]
\excheck 不能把三个电压的\emph{有效值}线性相加，
相量余弦项是关键。

\electricoscproblem{142}{173--174}
交流电源向并联的线圈和纯电阻
$R=25\,\Omega$ 供电，
总电流有效值为 $I=0.90$A，
线圈支路 $I_1=0.50$A，
电阻支路 $I_2=0.60$A。
求线圈内部热功率。
\exsolution
两支路电压相同，纯电阻支路电流与电压同相，
因此电源有效电压
$V=I_2R=15$V。
线圈电流与电压相位差为 $\varphi$，
总电流相量模满足
$I^2=I_1^2+I_2^2+2I_1I_2\cos\varphi$。
解出
$\cos\varphi=(I^2-I_1^2-I_2^2)/(2I_1I_2)$。
线圈发热
\[
P_1=VI_1\cos\varphi
=\frac R2\left(I^2-I_1^2-I_2^2\right).
\]
代数值：
\[
\boxed{P_1=\frac{25}{2}(0.90^2-0.50^2-0.60^2)
=2.5\,\mathrm W}.
\]
\excheck 无感支路的电流不能与线圈支路的电流模直接相加得到总电流。

\electricoscproblem{143}{174}
电容 $C=73\,\mathrm{\mu F}$ 与纯电阻
$R=100\,\Omega$ 并联，
通以角频率 $\omega=314\,\mathrm{s^{-1}}$ 的正弦交流电。
求整个回路的阻抗模。
\exsolution
用复数导纳
$\underline Y=1/R+i\omega C$，
模
$|Y|=\sqrt{1/R^2+\omega^2C^2}$，
因此
\[
\boxed{|Z|=\frac1{|Y|}
=\frac R{\sqrt{1+(\omega RC)^2}}}.
\]
代 $\omega RC=314(100)(73\times10^{-6})\simeq2.29$，
\[
\boxed{|Z|\simeq40\,\Omega}.
\]
\excheck 并联电容增加总导纳，故总阻抗模比单独电阻 $100\Omega$ 小。

\electricoscproblem{144}{174}
图4.35有三种交流\emph{并联}电路，
外加端子 A、B 的正弦电压为参考相位：
(a) 纯电阻 $R$ 与内阻为 $r$、自感为 $L$ 的线圈并联；
(b) 电容 $C$ 与内阻为 $R$、自感为 $L$ 的线圈并联；
(c) 纯电阻 $R$ 与串联 LC 支路并联。
各器件参数均满足\emph{总电流滞后}电压角度 $\varphi$。
给出各支路与总电流相量的图示构造。
\exsolution
以外加电压相量 $\underline V=V$ 为正实轴。
(a) 电阻电流
$\underline I_R=V/R$ 沿正实轴，
线圈电流
$\underline I_{Lr}=V/(r+i\omega L)$ 位于第四象限，
总电流
\[
\boxed{\underline I_0=\frac VR+\frac V{r+i\omega L}}.
\]
(b) 电容电流
$\underline I_C=i\omega CV$ 沿正虚轴，
线圈电流
$\underline I_{LR}=V/(R+i\omega L)$ 位于第四象限，
$\boxed{\underline I_0=\underline I_C+\underline I_{LR}}$；
题设总电流滞后要求线圈负虚部的模大于电容正虚部。
(c) 电阻电流在正实轴，
串联 LC 分支电流
$\underline I_{LC}=V/[i(\omega L-1/(\omega C))]$。
要使总电流滞后，必须
$\omega L>1/(\omega C)$，
从而 $\underline I_{LC}$ 沿负虚轴，
\[
\boxed{\underline I_0=
\frac VR-\frac{iV}{\omega L-1/(\omega C)}}.
\]
\excheck 三图都用矢量加法\emph{按头尾相接}得总电流；
(a)(b) 中线圈电流同时具有实、负虚分量，
(c) 理想 LC 支路电流纯无功。

\electricoscproblem{145}{174}
电容 $C=1.0\,\mathrm{\mu F}$ 与
具有有功电阻 $R=0.10\,\Omega$、
自感 $L=1.0\,\mathrm{mH}$ 的线圈并联，
接有效电压 $V=31$V 的交流源。
求：(a) 并联电流谐振角频率；
(b) 谐振时供电总电流、线圈支路电流和电容支路电流的有效值。
\exsolution
总导纳
$Y=i\omega C+1/(R+i\omega L)$。
电流谐振要求其虚部为零：
$\omega C-\omega L/(R^2+\omega^2L^2)=0$，
故
\[
\boxed{\omega_{\rm res}
=\sqrt{\frac1{LC}-\frac{R^2}{L^2}}
\simeq3.0\times10^4\,\mathrm{rad/s}}.
\]
此时
$R^2+\omega_{\rm res}^2L^2=L/C$，
总导纳实部
$G=R/(L/C)=RC/L$，
总供电电流
\[
\boxed{I_{\rm supply}=VRC/L
\simeq3.1\,\mathrm{mA}}.
\]
线圈支路
$\boxed{I_{\rm coil}
=V/\sqrt{L/C}=V\sqrt{C/L}\simeq0.98\,\mathrm A}$；
电容支路
$\boxed{I_C=\omega_{\rm res}CV\simeq0.98\,\mathrm A}$。
\excheck 两支路电流很大却近于反相，导致电源总电流仅几毫安；
书末各处将其约写为 $3$mA、$1$A。

\electricoscproblem{146}{174}
电容 $C$ 与有功电阻为 $R$、自感为 $L$ 的线圈并联，
交流源角频率为 $\omega$。
求总供电电流相对电源电压的相位差。
\exsolution
以源电压沿复平面实轴，
总导纳
\[
\underline Y=i\omega C+\frac1{R+i\omega L}
=\frac R{R^2+\omega^2L^2}
+i\left[\omega C-\frac{\omega L}{R^2+\omega^2L^2}\right].
\]
电流相位即 $\arg Y$，
因此
\[
\boxed{\tan\varphi=
\frac{\omega C(R^2+\omega^2L^2)-\omega L}{R}}.
\]
\excheck $\varphi>0$ 时总电流超前电压（容性）；
$\varphi<0$ 时滞后电压（感性）；$\varphi=0$ 为并联谐振。

\electricoscproblem{147}{174}
有功电阻为 $R$、自感为 $L$ 的线圈
与电容 $C$ 并联，
求交流角频率 $\omega$ 下整个回路的阻抗模。
\exsolution
由并联导纳
$Y=i\omega C+1/(R+i\omega L)$，
可写
\[
Y=\frac{1+i\omega C(R+i\omega L)}
{R+i\omega L}
=\frac{(1-\omega^2LC)+i\omega CR}
{R+i\omega L}.
\]
故
\[
\boxed{|Z|=\frac1{|Y|}
=\sqrt{\frac{R^2+\omega^2L^2}
{(1-\omega^2LC)^2+(\omega CR)^2}}}.
\]
\excheck $\omega\to0$ 时 $Z\to R$，
符合电容开路、线圈仅剩有功电阻的极限。

\electricoscproblem{148}{174}
细金属线环有电阻 $R$、自感 $L$，
在均匀外磁场中以恒定角速度 $\omega$ 转动，
外磁通随时间
$\Phi(t)=\Phi_0\cos\omega t$。
证明：(a) 感应电流的稳态规律；
(b) 维持转动所需外力平均机械功率。
\exsolution
法拉第电动势
$\mathcal E=-\dot\Phi=\omega\Phi_0\sin\omega t$。
回路
$L\dot I+RI=\mathcal E$。
其稳态解写为
\[
\boxed{I=I_m\sin(\omega t-\varphi)},\qquad
\boxed{I_m=\frac{\omega\Phi_0}
{\sqrt{R^2+\omega^2L^2}}},
\quad
\boxed{\tan\varphi=\frac{\omega L}{R}}.
\]
外力维持匀速转动所做的机械功在稳态全部补偿电阻损耗；
一周期平均
$\langle P_{\rm mech}\rangle
=R\langle I^2\rangle=RI_m^2/2$。
故
\[
\boxed{\langle P_{\rm mech}\rangle
=\frac{\omega^2\Phi_0^2R}{2(R^2+\omega^2L^2)}}.
\]
\excheck 线圈电感可在瞬时交换储能，但一周期首末磁场能相同，平均机械输入等于焦耳热。

\electricoscproblem{149}{175}
图4.36：木芯上沿轴线放两只线圈。
线圈1的电流
$I_1(t)=I_0\cos\omega t$；
线圈2短路，其有功电阻 $R$、自感 $L_2$，
互感 $M=M(x)$ 随两线圈间距 $x$ 改变。
求在稳态下两线圈间相互作用力沿 $x$ 的周期平均值。
\exsolution
线圈2的法拉第方程
$L_2\dot I_2+RI_2=-M(x)\dot I_1$，
以 $x$ 在一个振荡周期内不变近似。
复相量为
$\underline I_1=I_0$，
$\underline I_2=
-i\omega M I_0/(R+i\omega L_2)$，
所以
$\Re\underline I_2=
-\omega^2M L_2 I_0/(R^2+\omega^2L_2^2)$。
耦合磁通导致瞬时轴向力
$F_x=I_1I_2\,\mathrm dM/\mathrm dx$
（取正 $x$ 为两线圈拉开方向），
故周期平均
\[
\boxed{\langle F_x\rangle=
-\frac{\omega^2L_2I_0^2}{2(R^2+\omega^2L_2^2)}
\,M(x)\frac{\mathrm dM}{\mathrm dx}}.
\]
\excheck 对同绕向线圈，互感随距离增加通常减小（$M>0$、$M'<0$），所以
$\langle F_x\rangle>0$ 表示感生电流产生平均排斥力。
书后简答的符号取决于所选 $x$ 方向与互感有向约定，须一并指定。
% === END 24-electric-oscillations-138-149.tex ===
